% !TeX root = ../../main.tex
\chapter{算法分析}[Algorithm Analysis]\label{cw:sec:algorithms}
本节先给出每个时刻都有观测的基本算法。输入观测序列 $y_{1:T}=(y_1,\dots,y_T)$, 其中 $y_k$ 是时刻 $k$ 实际获得的观测值, $T$ 是总时间步数。已知状态转移模型 $f$, 过程噪声分布 $\mathbb P_{N_x}$ 与观测条件密度 $p_{Y_k\mid X_k}(y\mid x)$；收到 $y_k$ 后定义似然函数
\begin{equation}
 g_k(x):=p_{Y_k\mid X_k}(y_k\mid x).
\end{equation}
两种算法都从初始密度 $p(x_0)$ 采样, 每轮依次预测粒子, 用当前观测更新权重, 再生成供下一轮使用的等权粒子。假设每轮的粒子似然总和有限且严格为正。

\textbf{初始化与第一次迭代}\enspace 第一次观测到达时, 精确的预测与更新为
\begin{equation}\label{cw:eq:algorithm-first-bayes}
\begin{aligned}
&p(x_1)=\int p(x_1\mid x_0)p(x_0)\,\mathrm d x_0,\\
&p(x_1\mid y_1)=\frac{g_1(x_1)p(x_1)}{\int g_1(z)p(z)\,\mathrm d z}.
\end{aligned}
\end{equation}
对应的粒子近似是
\begin{equation}\label{cw:eq:algorithm-first-particles}
\begin{aligned}
&x_0^{+,(i)}\sim p(x_0),\;n_{x,1}^{(i)}\sim\mathbb P_{N_x},\;x_1^{-,(i)}=f(x_0^{+,(i)},n_{x,1}^{(i)}),\\
&w_1^{+,(i)}=\frac{g_1(x_1^{-,(i)})}{\sum_jg_1(x_1^{-,(j)})},\;
\hat{\mathbb P}_{N;X_1}^+=\sum_iw_1^{+,(i)}\delta_{x_1^{-,(i)}}\approx\mathbb P_{X_1\mid Y_1=y_1}.
\end{aligned}
\end{equation}
随后 BPF 按权重随机复制粒子, ETPF 通过最优输运取重心, 所得 $x_1^{+,(i)}$ 均作为第二轮预测的输入。

\section{自举粒子滤波}\label{cw:sec:alg-bpf}
基本 BPF 采用转移先验提议与独立重采样, 即 Gordon 等\cite{gordon1993novel}的预测与抽样更新。给定权重后, 每个新粒子独立选取一个旧粒子, 第 $i$ 个旧粒子被选中的概率为 $w_k^{+,(i)}$；复制次数因此服从多项分布。

算法~\ref{cw:alg:bpf} 的权重更新由 Bayes 公式得到；重采样只改变粒子的位置及复制次数；每轮总是从等权粒子开始（亦即，总是暴力重采样而不考虑有效粒子数）。

\textbf{间隔观测}\enspace 若只在 $\mathcal O\subseteq\{1,\dots,T\}$ 提供观测, 输入改为带时刻的观测记录 $\{(k,y_k):k\in\mathcal O\}$。在 $k\notin\mathcal O$ 时, 两种算法均只执行预测, 令 $x_k^{+,(i)}=x_k^{-,(i)}$, 输出 $\hat x_k=N^{-1}\sum_i x_k^{-,(i)}$, 然后进入下一时刻。每隔 $m$ 步观测对应 $\mathcal O=\{m,2m,\dots\}\cap\{1,\dots,T\}$。
% 先完成文字说明，再放置完整伪代码；避免不可拆分算法框提前截断正文页。
% !TeX root = ../../../main.tex
% 模板原生算法样式；计算步骤与迁入源稿一致。
\begin{algorithm}[H]
\caption{BPF：预测、似然加权与独立重采样}\label{cw:alg:bpf}
\KwIn{初始密度 $p(x_0)$, 粒子数 $N$, 模型 $f$, 噪声分布 $\mathbb P_{N_x}$, 观测条件密度, 观测序列 $y_{1:T}$}
\KwOut{后验均值估计 $\hat x_{1:T}$；每轮更新后的等权粒子 $\{x_k^{+,(j)}\}_{j=1}^N$}
独立采样 $x_0^{+,(i)}\sim p(x_0)$, 权重均为 $1/N$, $i=1,\dots,N$\;
\For{$k=1,\dots,T$}{
    独立采样 $n_{x,k}^{(i)}\sim\mathbb P_{N_x}$, 令 $x_k^{-,(i)}=f(x_{k-1}^{+,(i)},n_{x,k}^{(i)})$\;
    读取 $y_k$, 计算 $g_k(x_k^{-,(i)})=p_{Y_k\mid X_k}(y_k\mid x_k^{-,(i)})$\;
    $\displaystyle w_k^{+,(i)}=\frac{g_k(x_k^{-,(i)})}{\sum_{r=1}^N g_k(x_k^{-,(r)})}$, $i=1,\dots,N$\;
    输出后验均值估计 $\hat x_k=\sum_{i=1}^Nw_k^{+,(i)}x_k^{-,(i)}$\;
    给定 $\bm w_k^+$, 独立抽取 $A_1,\dots,A_N$, 其中 $\mathbb P(A_j=i\mid\bm w_k^+)=w_k^{+,(i)}$\;
    $x_k^{+,(j)}\gets x_k^{-,(A_j)}$, 新粒子权重均置为 $1/N$, $j=1,\dots,N$\;
}
\end{algorithm}


\section{集合输运粒子滤波}\label{cw:sec:alg-etpf}
Reich 的原始 ETPF\cite{reich2013nonparametric}共享 BPF 的预测与权重更新, 通过离散最优输运替换随机重采样。令 $\bm a=\bm w_k^+$, $\bm b=\mathbf1_N/N$, $C_{ij}=\|x_k^{-,(i)}-x_k^{-,(j)}\|^2$, 求解线性规划
\begin{equation}\label{cw:eq:algorithm-etpf-original}
\bm\Pi_k\in\operatorname*{argmin}_{\bm\Pi\in U(\bm a,\bm b)}\langle\bm\Pi,\bm C\rangle_F.
\end{equation}
目标槽位 $j$ 的参考位置为 $x_k^{-,(j)}$。由于列和为 $1/N$, $N\Pi_{k;ij}$ 是该槽位对应的条件权重, 新粒子取其加权平均。

% !TeX root = ../../../main.tex
% 模板原生算法样式；计算步骤与迁入源稿一致。
\begin{algorithm}[H]
\caption{ETPF：最优输运与重心变换}\label{cw:alg:etpf}
\KwIn{与算法~\ref{cw:alg:bpf} 相同的输入}
\KwOut{后验均值估计 $\hat x_{1:T}$；每轮变换后的等权粒子 $\{x_k^{+,(j)}\}_{j=1}^N$}
独立采样 $x_0^{+,(i)}\sim p(x_0)$, 权重均为 $1/N$, $i=1,\dots,N$\;
\For{$k=1,\dots,T$}{
    独立采样 $n_{x,k}^{(i)}\sim\mathbb P_{N_x}$, 令 $x_k^{-,(i)}=f(x_{k-1}^{+,(i)},n_{x,k}^{(i)})$\;
    读取 $y_k$, 计算 $g_k(x_k^{-,(i)})=p_{Y_k\mid X_k}(y_k\mid x_k^{-,(i)})$\;
    $\displaystyle w_k^{+,(i)}=\frac{g_k(x_k^{-,(i)})}{\sum_{r=1}^N g_k(x_k^{-,(r)})}$, $i=1,\dots,N$\;
    输出后验均值估计 $\hat x_k=\sum_{i=1}^Nw_k^{+,(i)}x_k^{-,(i)}$\;
    $\bm a\gets\bm w_k^+$, $\bm b\gets\mathbf1_N/N$, $C_{ij}\gets\|x_k^{-,(i)}-x_k^{-,(j)}\|^2$\;
    求解式~\eqref{cw:eq:algorithm-etpf-original}, 得到最优耦合 $\bm\Pi_k$\;
    $x_k^{+,(j)}\gets N\sum_{i=1}^N\Pi_{k;ij}x_k^{-,(i)}$, 权重均置为 $1/N$, $j=1,\dots,N$\;
}
\end{algorithm}


\textbf{熵正则化版本与基本 Sinkhorn 迭代}\enspace 使用式~\eqref{cw:eq:regularized-matrix} 的熵正则化问题\cite{cuturi2013sinkhorn}近似上述线性规划, 将算法~\ref{cw:alg:etpf} 中的耦合求解替换为算法~\ref{cw:alg:sinkhorn-basic}, 再执行相同的重心变换。固定 $\varepsilon>0$ 时得到的是正则化 ETPF, 与原始 ETPF 的输运目标有所不同。达到迭代上限时仍可返回近似耦合, 其边际误差由 $r$ 表示。

% !TeX root = ../../../main.tex
% 模板原生算法样式；计算步骤与迁入源稿一致。
\begin{algorithm}[H]
\caption{基本 Sinkhorn 迭代：求解正则化耦合}\label{cw:alg:sinkhorn-basic}
\KwIn{边际 $\bm a,\bm b$, 代价矩阵 $\bm C$, $\varepsilon>0$, 容差 $\tau>0$, 最大迭代次数 $L_{\max}\geq1$}
\KwOut{近似耦合 $\bm\Pi$、边际残差 $r$ 及收敛标志 $r\leq\tau$}
$K_{ij}\gets\exp(-C_{ij}/\varepsilon)$, $\bm v\gets\mathbf1_N$\;
\For{$\ell=1,\dots,L_{\max}$}{
    $\bm u\gets\bm a\oslash(\bm K\bm v)$\;
    $\bm v\gets\bm b\oslash(\bm K^{\mathrm T}\bm u)$\;
    $\bm\Pi\gets\operatorname{diag}(\bm u)\bm K\operatorname{diag}(\bm v)$\;
    $r\gets\max\{\|\bm\Pi\mathbf1_N-\bm a\|_1,\|\bm\Pi^{\mathrm T}\mathbf1_N-\bm b\|_1\}$\;
    \If{$r\leq\tau$}{退出迭代\;}
}
返回 $\bm\Pi$, 残差 $r$ 与是否满足 $r\leq\tau$\;
\end{algorithm}


% 保持算法所属小节与误差分析的阅读顺序。
% !TeX root = ../../../main.tex
\section{均方误差的四项分解}\label{cw:sec:algorithm-errors}
令 $\mathcal H=L^2(\Omega,\mathcal F,\mathbb P;\mathbb R^{n_x})$，概率空间包含真实系统与算法随机性。
记 $\mathcal Y_k\subseteq\mathcal H$ 为所有关于 $\sigma(Y_{1:k})$ 可测的随机变量构成的闭子空间，
$\mathcal P_{\mathcal Y_k}X_k$ 是利用全部历史观测的最优均方估计。
本节记 $\hat X_k=X_k^+\mathbf1_N/N$ 为重采样后均值；算法~\ref{cw:alg:bpf} 与~\ref{cw:alg:etpf} 的在线输出 $X_k^-\bm w_k^+$ 是重采样前的加权均值。

\begin{theorem}[四项误差分解]\label{cw:thm:four-errors}
假设所涉及随机变量平方可积，算法随机源与真实系统独立。
对每步使用新的随机数进行多项重采样的 BPF，或满足精确边际约束的 ETPF，有
\begin{equation}\label{cw:eq:four-errors}
\begin{aligned}
&\|X_k-\hat X_k\|_{\mathcal H}^2\\
= &\underbrace{\|\mathcal P_{\mathcal Y_k}^{\perp}X_k\|_{\mathcal H}^2}_{\text{不可约 Bayes 误差}}
+\underbrace{\|\mathcal P_{\mathcal Y_k}(X_k-X_k^-\bm w_k^+)\|_{\mathcal H}^2}_{\text{算法偏差平方}}\\
&+\underbrace{\|\mathcal P_{\mathcal Y_k}^{\perp}(X_k^-\bm w_k^+)\|_{\mathcal H}^2}_{\text{算法随机波动}}
+\underbrace{\|X_k^-\bm w_k^+-\hat X_k\|_{\mathcal H}^2}_{\text{当前重采样扰动}}.
\end{aligned}
\end{equation}
\end{theorem}
\begin{proof}
令 $\mathcal Z_k$ 为截至当前的观测 $Y_{1:k}$ 与本步重采样前全部算法随机源生成的 $L^2$ 子空间；加权均值 $X_k^-\bm w_k^+$ 对相应的 $\sigma$ 代数可测，属于该子空间。
随机源独立性给出 $\mathcal P_{\mathcal Z_k}X_k=\mathcal P_{\mathcal Y_k}X_k$；
重采样的条件无偏性给出 $\mathcal P_{\mathcal Z_k}\hat X_k=X_k^-\bm w_k^+$，ETPF 则逐次精确相等。
当前重采样使用新的独立随机数，故其扰动与 $X_k$ 及 $\mathcal Z_k$ 正交。
\end{proof}

多项重采样 BPF 的最后一项可以直接计算：
\begin{equation}\label{cw:eq:bpf-resampling-error}
\|X_k^-\bm w_k^+-\hat X_k\|_{\mathcal H}^2
=\frac1N\sum_{i=1}^N\left\|\sqrt{w_k^{+,(i)}}\bigl(x_k^{-,(i)}-X_k^-\bm w_k^+\bigr)\right\|_{\mathcal H}^2.
\end{equation}
ETPF 由 $\bm\Pi_k\mathbf1_N=\bm w_k^+$ 得到
$X_k^+\mathbf1_N/N=X_k^-\bm\Pi_k\mathbf1_N=X_k^-\bm w_k^+$，因此最后一项为零。

不过，质心守恒并不能消除算法偏差与随机波动；有限次 Sinkhorn 迭代的边际残差也可能破坏精确守恒。因而实际数值结果会与理想计算有所不同。

\begin{theorem}[零过程噪声下的粒子贫化]\label{cw:thm:bpf-impoverishment}
固定 $N\ge2$，预测为确定性映射，每步权重有效且使用新的随机数进行多项重采样。
固定合法观测记录，记 $D_k:=\#\{x_k^{+,(i)}:i=1..N\}$，则
\begin{equation}\label{cw:eq:particle-count-decay}
1\le\mathbb E[D_k]\le1+(N-1)(1-N^{1-N})^k,\qquad
\lim_{k\to\infty}\mathbb E[D_k]=1.
\end{equation}
\end{theorem}
\begin{proof}
确定性预测与复制使 $D_k$ 不增，单点状态为吸收态。
每次全部选中同一祖先的条件概率为 $\sum_i(w_k^{+,(i)})^N\ge N^{1-N}$。
因此 $\mathbb P(D_k>1)\le(1-N^{1-N})^k$，再由 $D_k-1\le(N-1)\mathbb I_{\{D_k>1\}}$ 得证。
\end{proof}
更多关于粒子谱系合并的相关结果见
Jacob、Murray 与 Rubenthaler 的研究\cite{jacob2015path}；式~\eqref{cw:eq:particle-count-decay} 是上面的直接估计。

